\pagestyle{plain}
\chapter*{Agradecimientos}

La vida está repleta de casualidades. Si no fuera por todas las que se dieron, no estaría hoy aquí escribiendo esto, no sería quien soy hoy sin todo lo que me ha sucedido hasta ahora, el barco de Teseo no es el original, sino el que queda tras tanto cambio, el del día a día.

Quiero dar las gracias a todas esas casualidades que pasaron por mi vida y que hicieron de esta algo maravilloso. Me gustaría ir por orden cronológico, así que comenzaremos por quien lo inició todo. Me gustaría dar gracias a mis padres, hermana y cuñado, que me levantaron de mis mayores caídas y celebraron mis mayores alegrías, quienes, aún sin tener nada, me dieron todo y quienes, aún derrumbados, supieron levantarme una sonrisa y sonreír ellos mismos, que sabemos que hay veces que es muy difícil. Me gustaría, por supuesto, dar gracias al resto de mi familia, que se que siempre estarán ahí para lo que necesite. Me gustaría dar las gracias a mi mejor amigo y su familia, que nos entendemos desde que no podíamos ni hablar, que los considero de mi propia sangre aunque de esta no tengan. Me gustaría dar las gracias a todos los profesores que pasaron por mi vida. Aunque no todos fueran perfectos o no congeniáramos, me formaron, que es mucho más de lo que puede decir mucha gente, de todos aprendí aunque fuera algo. Concretamente, Montse, mi profesora de religión: gracias por enseñarme a "observar el remolino", creo que, gracias a eso, disfruto un poco más de las tretas de la vida. Quiero dar gracias a mis amigos, esa segunda familia que me lleva acompañando desde que tengo uso de razón y que me demostraron lo que es la confianza, tanto en el resto como en uno mismo, sin esos deberes a última hora mal copiados, sin esas conversaciones profundas a última hora de la noche y sin esas disculpas con la cabeza baja cuando se que lo hice mal, no sería el mismo que soy ahora. Siempre estaré agradecido por tener la oportunidad de compartir mi vida con ellos. Por último, quería dar las gracias a todos los que ya no están conmigo. La vida no es perfecta, lo que implica que a veces toca perder. Doy las gracias por esos "te quiero" aunque ya no los pueda escuchar, doy las gracias por esos "estoy orgullosa de ti", aunque ya no los pueda escuchar y doy las gracias por esas lágrimas que derramé, porque eso significa que sentí verdaderos sentimientos de afecto por vosotros.

En definitiva, doy las gracias por poder vivir y por poder ser quien soy. Muchas gracias.